\subsection{Polynomial Modules from Operators}

Throughout this subsection, $\mathcal{V}$ is finite-dimensional over $\mathbb{F}$ unless explicitly stated otherwise.

\begin{perspective} \label{per:operator-as-polynomial-action} \hlt{An Operator as a Polynomial Action}\\
Fix $T\in\LL(\mathcal{V})$. Evaluation at $T$ defines an $\mathbb{F}[z]$-action on $\mathcal{V}$ by $p\cdot\mathbf{v}\defeq p(T)\mathbf{v}$.\\
Thus the orbit $\mathbf{v},T\mathbf{v},T^2\mathbf{v},\ldots$ is encoded by the cyclic subspace $\{p(T)\mathbf{v}:p\in\mathbb{F}[z]\}=\operatorname{span}\{T^j\mathbf{v}:j\geq0\}$.
\end{perspective}

\begin{definition} \label{def:operator-module-and-matrix-module} \hlt{The Module Attached to an Operator}\\
Let $T\in\LL(\mathcal{V})$. The \hlt{operator module} attached to $T$ is the $\mathbb{F}[z]$-module $\mathcal{V}_T$ with underlying additive group $\mathcal{V}$ and action $p(z)\cdot\mathbf{v}\defeq p(T)\mathbf{v}$.
For $A\in\operatorname{Mat}(n,n,\mathbb{F})$, define $(\mathbb{F}^n)_A$ analogously by $p(z)\cdot\mathbf{x}\defeq p(A)\mathbf{x}$.
\end{definition}

\begin{proposition} \label{prop:operator-modules-are-fg-torsion} \hlt{Operator Modules Are Finitely Generated Torsion Modules}\\
Let $0<\dim\mathcal{V}<\infty$ and $T\in\LL(\mathcal{V})$. Then $\mathcal{V}_T$ is a finitely generated torsion $\mathbb{F}[z]$-module.
\begin{proof}
Let $(\mathbf{v}_1,\ldots,\mathbf{v}_n)$ be an $\mathbb{F}$-basis of $\mathcal{V}$. Since $\mathbb{F}\subseteq\mathbb{F}[z]$, this basis generates $\mathcal{V}_T$.\\
By Proposition \ref{prop:finite-dimensional-operators-have-annihilating-polynomials}, there exists $0\neq p\in\mathbb{F}[z]$ with $p(T)=0$. Thus $p\cdot\mathbf{v}=p(T)\mathbf{v}=\mathbf{0}$ for all $\mathbf{v}\in\mathcal{V}$, so $\mathcal{V}_T$ is torsion. Finally, $\mathbb{F}[z]$ is a PID by Theorem \ref{thm:integers-and-polynomial-ring-are-pids}.\\
\hlt{Structural Commentary}: Finite-dimensionality gives a global polynomial relation for the operator.\\
\hlt{Key Insight}: Rational canonical form is the coordinate form of the structure theorem for $\mathcal{V}_T$.
\end{proof}
\end{proposition}

\begin{theorem} \label{thm:operator-submodules-invariant-subspaces} \hlt{Submodules Are Invariant Subspaces}\\
Let $T\in\LL(\mathcal{V})$ and $\mathcal{S}\subseteq\mathcal{V}$. Then $\mathcal{S}\leq_{\mathbb{F}[z]}\mathcal{V}_T$ if and only if $\mathcal{S}\leq_{\mathbb{F}}\mathcal{V}$ and $T\mathcal{S}\subseteq\mathcal{S}$.
\begin{proof}
\begin{enumerate}[label=\roman*.]
\setlength{\itemsep}{0pt}
\item[$(\Rightarrow)$] If $\mathcal{S}\leq_{\mathbb{F}[z]}\mathcal{V}_T$, then $\mathbb{F}\subseteq\mathbb{F}[z]$ gives $\mathcal{S}\leq_{\mathbb{F}}\mathcal{V}$. Since $z\cdot\mathbf{s}=T\mathbf{s}$ for $\mathbf{s}\in\mathcal{S}$, we have $T\mathcal{S}\subseteq\mathcal{S}$.
\item[$(\Leftarrow)$] If $\mathcal{S}\leq_{\mathbb{F}}\mathcal{V}$ and $T\mathcal{S}\subseteq\mathcal{S}$, then $T^j\mathcal{S}\subseteq\mathcal{S}$ for all $j\geq0$.\\
Hence for $p(z)=\sum_ja_jz^j$ and $\mathbf{s}\in\mathcal{S}$, $p\cdot\mathbf{s}=p(T)\mathbf{s}=\sum_ja_jT^j\mathbf{s}\in\mathcal{S}$, so $\mathcal{S}\leq_{\mathbb{F}[z]}\mathcal{V}_T$.
\end{enumerate}
\hlt{Structural Commentary}: The single scalar $z$ remembers the operator.\\
\hlt{Key Insight}: Invariance under $T$ is exactly closure under all polynomial scalars.
\end{proof}
\end{theorem}

\begin{theorem} \label{thm:operator-similarity-iff-module-isomorphism} \hlt{Similarity Is Module Isomorphism}\\
Let $T,S\in\LL(\mathcal{V})$. Then $\mathcal{V}_T\cong_{\mathbb{F}[z]}\mathcal{V}_S\Longleftrightarrow T\sim S$. More precisely, for $\Phi\in\operatorname{GL}(\mathcal{V})$, $\Phi:\mathcal{V}_T\to\mathcal{V}_S$ is $\mathbb{F}[z]$-linear if and only if $\Phi T=S\Phi$, equivalently $S=\Phi T\Phi^{-1}$.
\begin{proof}
\begin{enumerate}[label=\roman*.]
\setlength{\itemsep}{0pt}
\item[$(\Rightarrow)$] If $\Phi:\mathcal{V}_T\to\mathcal{V}_S$ is an $\mathbb{F}[z]$-isomorphism, then $\Phi\in\operatorname{GL}(\mathcal{V})$, $\Phi(T\mathbf{v})=\Phi(z\cdot\mathbf{v})=z\cdot\Phi(\mathbf{v})=S\Phi(\mathbf{v})$ $\forall \mathbf{v}\in\mathcal{V}$.\\
Hence $\Phi T=S\Phi$, so $S=\Phi T\Phi^{-1}$.
\item[$(\Leftarrow)$] If $\Phi\in\operatorname{GL}(\mathcal{V})$ and $S=\Phi T\Phi^{-1}$, then $\Phi T=S\Phi$, hence $\Phi T^j=S^j\Phi$ for all $j\geq0$.\\
Thus for $p(z)=\sum_ja_jz^j$, $\Phi(p(T)\mathbf{v})=p(S)\Phi(\mathbf{v})$, so $\Phi:\mathcal{V}_T\to\mathcal{V}_S$ is an $\mathbb{F}[z]$-isomorphism.
\end{enumerate}
\hlt{Structural Commentary}: A module isomorphism is precisely an intertwiner for the action of $z$.\\
\hlt{Key Insight}: Similarity is isomorphism of the associated $\mathbb{F}[z]$-modules.
\end{proof}
\end{theorem}

\begin{corollary} \label{cor:matrix-similarity-iff-matrix-modules-isomorphic} \hlt{Matrix Similarity Is Matrix-Module Isomorphism}\\
Let $A,B\in\operatorname{Mat}(n,n,\mathbb{F})$. Then $A\sim B\Longleftrightarrow(\mathbb{F}^n)_A\cong_{\mathbb{F}[z]}(\mathbb{F}^n)_B$. More precisely, for $X\in\operatorname{GL}_n(\mathbb{F})$, $X:(\mathbb{F}^n)_A\to(\mathbb{F}^n)_B$ is $\mathbb{F}[z]$-linear if and only if $XA=BX$, equivalently $B=XAX^{-1}$.
\begin{proof}
\begin{enumerate}[label=\roman*.]
\setlength{\itemsep}{0pt}
\item[$(\Rightarrow)$] If $X:(\mathbb{F}^n)_A\to(\mathbb{F}^n)_B$ is an $\mathbb{F}[z]$-isomorphism, then $X\in\operatorname{GL}_n(\mathbb{F})$ and $X(A\mathbf{x})=X(z\cdot\mathbf{x})=z\cdot X\mathbf{x}=BX\mathbf{x}$ for all $\mathbf{x}\in\mathbb{F}^n$. Hence $XA=BX$, so $B=XAX^{-1}$.
\item[$(\Leftarrow)$] If $X\in\operatorname{GL}_n(\mathbb{F})$ and $B=XAX^{-1}$, then $XA=BX$, hence $XA^j=B^jX$ for all $j\geq0$. Thus $Xp(A)=p(B)X$ for all $p\in\mathbb{F}[z]$, so $X:(\mathbb{F}^n)_A\to(\mathbb{F}^n)_B$ is an $\mathbb{F}[z]$-isomorphism.
\end{enumerate}
\hlt{Structural Commentary}: Similarity is invertible module intertwining.\\
\hlt{Key Insight}: The matrix is remembered by the scalar action of $z$.
\end{proof}
\end{corollary}

\begin{definition} \label{def:operator-module-annihilator-minimal-polynomial} \hlt{Module Annihilator and Minimal Polynomial}
\begin{enumerate}[label=(\roman*)]
\setlength{\itemsep}{0pt}
\item The \hlt{annihilator} of $\mathcal{V}_T$ is $\operatorname{Ann}_{\mathbb{F}[z]}(\mathcal{V}_T)\defeq\{p\in\mathbb{F}[z]:p(T)\mathcal{V}=\{\mathbf{0}\}\}$.
\item \hlt{Minimal polynomial} $\mu_T$ is unique monic generator of this nonzero principal ideal, so $\operatorname{Ann}_{\mathbb{F}[z]}(\mathcal{V}_T)=\langle\mu_T\rangle$.
\item For $A\in\operatorname{Mat}(n,n,\mathbb{F})$, define $\mu_A\defeq\mu_{T_A}$, where $T_A(\mathbf{x})=A\mathbf{x}$.
\end{enumerate}
\end{definition}

\begin{proposition} \label{prop:minimal-polynomial-module-order} \hlt{Minimal Polynomial as Module Order}\\
Let $T\in\LL(\mathcal{V})$. Then $\mu_T$ is the unique monic polynomial of minimal degree in $\{p\in\mathbb{F}[z]:p(T)=0\}$. If $S\sim T$, then $\mu_S=\mu_T$; hence similar matrices have the same minimal polynomial.
\begin{proof}
By Definition \ref{def:operator-module-annihilator-minimal-polynomial}, $\operatorname{Ann}_{\mathbb{F}[z]}(\mathcal{V}_T)=\langle\mu_T\rangle$ and $p\in\operatorname{Ann}_{\mathbb{F}[z]}(\mathcal{V}_T)$ iff $p(T)=0$. Thus $\mu_T$ is the unique monic annihilator of minimal degree. If $S=\Phi T\Phi^{-1}$, then $p(S)=\Phi p(T)\Phi^{-1}$ for all $p\in\mathbb{F}[z]$, so $\operatorname{Ann}_{\mathbb{F}[z]}(\mathcal{V}_S)=\operatorname{Ann}_{\mathbb{F}[z]}(\mathcal{V}_T)$ and $\mu_S=\mu_T$. The matrix case follows by applying this to the associated operators.\\
\hlt{Structural Commentary}: The minimal polynomial is the order of the whole operator module.\\
\hlt{Key Insight}: Similarity preserves the module and therefore preserves its annihilator ideal.
\end{proof}
\end{proposition}

\begin{definition} \label{def:cyclic-subspace-and-vector-annihilator} \hlt{Cyclic Subspace and Vector Annihilator}\\
Let $T\in\LL(\mathcal{V})$ and $\mathbf{v}\in\mathcal{V}$.
\begin{enumerate}[label=(\roman*)]
\setlength{\itemsep}{0pt}
\item The \hlt{$T$-cyclic subspace generated by $\mathbf{v}$} is $\mathcal{C}_T(\mathbf{v})\defeq\operatorname{span}\{T^j\mathbf{v}:j\geq0\}$
\item The vector $\mathbf{v}$ is a \hlt{cyclic vector} for $T$ if $\mathcal{C}_T(\mathbf{v})=\mathcal{V}$.
\item The \hlt{annihilator of $\mathbf{v}$ for $T$} is $\operatorname{ann}_T(\mathbf{v})\defeq\{p\in\mathbb{F}[z]:p(T)\mathbf{v}=\mathbf{0}\}$
\end{enumerate}
\end{definition}

\begin{proposition} \label{prop:cyclic-subspace-smallest-invariant-subspace} \hlt{Cyclic Subspaces Are Smallest Invariant Hosts}\\
Let $T\in\LL(\mathcal{V})$ and $\mathbf{v}\in\mathcal{V}$. Then $\mathcal{C}_T(\mathbf{v})=\{p(T)\mathbf{v}:p\in\mathbb{F}[z]\}$.\\
Moreover, $\mathcal{C}_T(\mathbf{v})$ is the smallest $T$-invariant subspace containing $\mathbf{v}$, and $\operatorname{ann}_T(\mathbf{v})$ is an ideal of $\mathbb{F}[z]$.
\begin{proof} Equality follows as $p(T)\mathbf{v}$ is a finite linear combination of the orbit vectors $T^j\mathbf{v}$, and every such finite linear combination equals $p(T)\mathbf{v}$ for some $p\in\mathbb{F}[z]$. Also $T(p(T)\mathbf{v})=(zp)(T)\mathbf{v}$, so $\mathcal{C}_T(\mathbf{v})$ is $T$-invariant.\\
If $\mathcal{U}$ is $T$-invariant and contains $\mathbf{v}$, then $T^j\mathbf{v}\in\mathcal{U}$ for every $j\geq0$, hence $\mathcal{C}_T(\mathbf{v})\subseteq\mathcal{U}$. Finally, if $p,q\in\operatorname{ann}_T(\mathbf{v})$ and $r\in\mathbb{F}[z]$, then $(p+q)(T)\mathbf{v}=\mathbf{0}$ and $(rp)(T)\mathbf{v}=r(T)p(T)\mathbf{v}=\mathbf{0}$ by Proposition \ref{prop:operator-polynomial-products}. $\operatorname{ann}_T(\mathbf{v})$ is an ideal.\\
\hlt{Structural Commentary}: A cyclic subspace is the span of the orbit generated by repeatedly applying $T$.\\
\hlt{Key Insight}: Invariance is exactly what lets the whole orbit stay inside the same host subspace.
\end{proof}
\end{proposition}

\begin{figure}[H]
  \centering
  \begin{tikzpicture}[
  >=Stealth,
  every node/.style={font=\small},
  orbit/.style={draw=blue!60!black, fill=blue!6, rounded corners=2pt,
    minimum width=1.15cm, minimum height=0.58cm, align=center, inner sep=4pt},
  subspace/.style={draw=dartmouthgreen!65!black, fill=dartmouthgreen!8, rounded corners=2pt,
    minimum width=4.7cm, minimum height=0.72cm, align=center, inner sep=6pt},
  arr/.style={->, line width=0.7pt, draw=black!70}
]
  \node[orbit] (v0) at (-2.45,0) {$\mathbf{v}$};
  \node[orbit] (v1) at (-0.85,0) {$T\mathbf{v}$};
  \node[orbit] (v2) at (0.85,0) {$T^2\mathbf{v}$};
  \node[orbit, minimum width=0.8cm] (dots) at (2.35,0) {$\cdots$};

  \draw[arr] (v0) -- node[above] {$T$} (v1);
  \draw[arr] (v1) -- node[above] {$T$} (v2);
  \draw[arr] (v2) -- node[above] {$T$} (dots);

  \node[subspace] (C) at (0,-1.35)
    {$\mathcal{C}_T(\mathbf{v})=\operatorname{span}\{T^j\mathbf{v}:j\geq0\}$};

  \draw[arr] (0,-0.38) -- node[right] {$\operatorname{span}$} (C.north);
\end{tikzpicture}

  \caption{The $T$-orbit of $\mathbf{v}$ spans $\mathcal{C}_T(\mathbf{v})$.}
  \label{fig:operator-cyclic-subspace-polynomial-action}
\end{figure}

\begin{definition} \label{def:vector-order-under-operator} \hlt{Order of a Vector under an Operator}\\
Let $\mathcal{V}$ be finite-dimensional, let $T\in\LL(\mathcal{V})$, and let $\mathbf{v}\neq\mathbf{0}$. The \hlt{order of $\mathbf{v}$ under $T$} is the unique monic polynomial $m_{\mathbf{v}}\in\operatorname{ann}_T(\mathbf{v})$ of least degree; equivalently, $\operatorname{ann}_T(\mathbf{v})=(m_{\mathbf{v}})$
\end{definition}

\begin{definition} \label{def:companion-matrix} \hlt{Companion Matrix}\\
Let $f(z)=z^n+c_{n-1}z^{n-1}+\cdots+c_1z+c_0$ be monic. For $n\geq2$, the \hlt{companion matrix} of $f$ is
\begin{align}
C_f\defeq
\begin{bmatrix}
0&0&\cdots&0&-c_0\\
1&0&\cdots&0&-c_1\\
0&1&\cdots&0&-c_2\\
\vdots&\vdots&\ddots&\vdots&\vdots\\
0&0&\cdots&1&-c_{n-1}
\end{bmatrix}.
\nonumber
\end{align}
For $n=1$, set $C_f=[-c_0]$.
\end{definition}

\begin{theorem} \label{thm:companion-matrix-has-prescribed-minimal-polynomial} \hlt{Companion Matrix Realises a Polynomial}\\
Let $f(z)=z^n+\sum_{i=0}^{n-1}c_iz^i$ be monic and let $C_f$ be its companion matrix. Then $p_{C_f}=\mu_{C_f}=f$.
\begin{proof}
Let $(\mathbf{e}_1,\ldots,\mathbf{e}_n)$ be the standard basis.\\
For $0\leq j<n$, $C_f^j\mathbf{e}_1=\mathbf{e}_{j+1}$, and last column gives $C_f^n\mathbf{e}_1=-\sum\limits_{i=0}^{n-1}c_iC_f^i\mathbf{e}_1$. Hence $f(C_f)\mathbf{e}_1=0$.\\
Since $\mathbf{e}_j=C_f^{j-1}\mathbf{e}_1$ and $f(C_f)C_f=C_ff(C_f)$, $f(C_f)\mathbf{e}_j=0$ $\forall j$, so $f(C_f)=0$.\\
If $0\neq g(z)=\sum\limits_{i=0}^m b_iz^i$ with $m<n$, then $g(C_f)\mathbf{e}_1=\sum\limits_{i=0}^m b_i\mathbf{e}_{i+1}\neq\mathbf{0}$.\\
Thus no nonzero polynomial of degree $<n$ annihilates $C_f$, so $\deg\mu_{C_f}=n$.\\
Since $\mu_{C_f}\mid f$ and both are monic of degree $n$, $\mu_{C_f}=f$.\\
By Cayley--Hamilton, $\mu_{C_f}\mid p_{C_f}$; as $p_{C_f}$ and $f$ are monic of degree $n$, $p_{C_f}=f$.\\
\hlt{Structural Commentary}: The cyclic vector $\mathbf{e}_1$ generates the basis, and last column imposes relation $f(C_f)=0$.\\
\hlt{Key Insight}: The companion matrix makes $f$ the first possible annihilating relation.
\end{proof}
\end{theorem}

\begin{theorem} \label{thm:cyclic-subspace-quotient-and-companion} \label{prop:cyclic-basis-gives-companion-matrix} \hlt{Cyclic Subspaces as Polynomial Quotients}\\
Let $\dim\mathcal{V}<\infty$, $T\in\LL(\mathcal{V})$, $0\neq\mathbf{v}\in\mathcal{V}$, and let $m_{\mathbf{v}}(z)=z^d+\sum\limits_{i=0}^{d-1}c_iz^i$ be the order of $\mathbf{v}$.\\
Define $\Phi_{\mathbf{v}}:\mathbb{F}[z]\to\mathcal{C}_T(\mathbf{v})$ by $\Phi_{\mathbf{v}}(p)=p(T)\mathbf{v}$ and $\mathcal{B}_{\mathbf{v}}\defeq(\mathbf{v},T\mathbf{v},\ldots,T^{d-1}\mathbf{v})$. Then
\begin{enumerate}[label=(\roman*)]
\setlength{\itemsep}{0pt}
\item $\operatorname{im}\Phi_{\mathbf{v}}=\mathcal{C}_T(\mathbf{v})$, $\ker\Phi_{\mathbf{v}}=\langle m_{\mathbf{v}}\rangle$, and $\mathcal{C}_T(\mathbf{v})\cong\mathbb{F}[z]/\langle m_{\mathbf{v}}\rangle$;
\item $\mathcal{B}_{\mathbf{v}}$ is a basis of $\mathcal{C}_T(\mathbf{v})$;
\item $[T|_{\mathcal{C}_T(\mathbf{v})}]_{\mathcal{B}_{\mathbf{v}}}=C_{m_{\mathbf{v}}}$.
\end{enumerate}
\begin{proof}
\begin{enumerate}[label=(\roman*)]
\setlength{\itemsep}{0pt}
\item By Proposition \ref{prop:cyclic-subspace-smallest-invariant-subspace}, $\operatorname{im}\Phi_{\mathbf{v}}=\mathcal{C}_T(\mathbf{v})$. $\ker\Phi_{\mathbf{v}}=\operatorname{ann}_T(\mathbf{v})=\langle m_{\mathbf{v}}\rangle$, by Theorem \ref{thm:quotient-by-kernel-is-range} $\mathcal{C}_T(\mathbf{v})\cong\mathbb{F}[z]/\langle m_{\mathbf{v}}\rangle$.
\item If $\sum\limits_{j=0}^{d-1}a_jT^j\mathbf{v}=\mathbf{0}$, then $q(z)=\sum\limits_{j=0}^{d-1}a_jz^j\in\operatorname{ann}_T(\mathbf{v})$ and $\deg q<d$, hence $q=0$; thus $\mathcal{B}_{\mathbf{v}}$ is independent.\\
For $p=hm_{\mathbf{v}}+r$ with $\deg r<d$ or $r=0$, $p(T)\mathbf{v}=r(T)\mathbf{v}\in\operatorname{span}\mathcal{B}_{\mathbf{v}}$, so $\mathcal{B}_{\mathbf{v}}$ spans.
\item For $0\leq j<d-1$, $T(T^j\mathbf{v})=T^{j+1}\mathbf{v}$. Since $m_{\mathbf{v}}(T)\mathbf{v}=\mathbf{0}$, $T^d\mathbf{v}=-\sum\limits_{i=0}^{d-1}c_iT^i\mathbf{v}$, last-column relation of $C_{m_{\mathbf{v}}}$.
\end{enumerate}
\hlt{Structural Commentary}: The first relation cuts out the quotient and companion block.\\
\hlt{Key Insight}: One orbit carries the whole cyclic subspace.
\end{proof}
\end{theorem}

\begin{perspective} \label{per:module-operator-dictionary} \hlt{The Translation Dictionary}\\
The module vocabulary for $\mathcal{V}_T$ and the vector-space vocabulary for $T$ match as follows:
\begin{center}
\begin{tabular}{c|c}
$\mathbb{F}[z]$-module $\mathcal{V}_T$ & $\mathbb{F}$-vector space $\mathcal{V}$\\
\hline
$p(z)\cdot\mathbf{v}$ & $p(T)\mathbf{v}$\\
submodule & $T$-invariant subspace\\
$\operatorname{Ann}_{\mathbb{F}[z]}(\mathcal{V}_T)=\langle\mu_T\rangle$ & minimal polynomial of $T$\\
$\langle\!\langle\mathbf{v}\rangle\!\rangle_{\mathbb{F}[z]}$ & $\mathcal{C}_T(\mathbf{v})$\\
$\mathbb{F}[z]/\langle m_{\mathbf{v}}\rangle$ & cyclic block with companion matrix $C_{m_{\mathbf{v}}}$
\end{tabular}
\end{center}
\hlt{Key Insight}: Once $z$ is interpreted as $T$, module structure becomes operator structure.
\end{perspective}

\begin{method} \label{meth:operator-module-to-spectral-form} \hlt{Moving Between Module and Spectral Data}\\
For $T\in\LL(\mathcal{V})$, translate the $\mathbb{F}[z]$-module data of $\mathcal{V}_T$ into spectral data as follows.
\begin{enumerate}[label=(\roman*)]
\setlength{\itemsep}{0pt}
\item Factor $\mu_T=\prod_{i=1}^s p_i^{e_i}$ over $\mathbb{F}$; the primary components are $\ker p_i(T)^{e_i}$.
\item Decompose each $\ker p_i(T)^{e_i}$ into cyclic summands; their orders $p_i^{e_{i,j}}$ are the elementary divisors.
\item Over $\mathbb{F}$, each elementary divisor $p_i^{e_{i,j}}$ gives a companion block $C_{p_i^{e_{i,j}}}$ in rational canonical form.
\item Over splitting field for $p_i$, same primary cyclic summands refine into Jordan blocks (Perspective \ref{per:elementary-divisors-and-jordan-blocks}).
\item $T$ is diagonalizable exactly when every elementary divisor is linear, i.e., Corollary \ref{cor:diagonal-rational-canonical-form-criterion} applies.
\end{enumerate}
\hlt{Key Insight}: Rational and Jordan forms encode the same elementary-divisor data over different fields.
\end{method}
